\section*{Supplementary Methods S1}
\providecommand{\suppmethodsection}[1]{\subsection*{#1}}

\suppmethodsection{Scope and evidence boundary}

This supplement describes source-data governance, orthogroup construction, rank-based programme discovery, validation and post hoc robustness analyses for PlantHOP. No predictive model was trained. Source H5AD and RDS objects were read without modification.

All labels are author-provided or previously governed source annotations. No label was inferred from expression during preparation of Tables S2 and S3. Orthogroup recurrence is treated as an association and portability device, not as proof of conserved molecular function, one-to-one orthology, causal marker activity, or paralog substitution. Missing values remain blank with explicit gap fields; they were never replaced by zero.

\suppmethodsection{Computing environment and recorded software}

Metadata processing used Python 3.12.14 with anndata 0.13.3.post0, NumPy 2.5.3, pandas 3.0.5 and SciPy 1.18.0. H5AD metadata were read in backed mode. Nested-label analyses used R 4.4.3; the retained environment currently contains Seurat 5.5.1, although the original run did not serialize every package version.

Orthogroups were generated with OrthoFinder 3.1.5 using DIAMOND sequence search and the recorded MSA/FAMSA/FastTree workflow. The seven-species run used 12 threads, contained 183,178 representative proteins, produced 18,383 orthogroups, and assigned 164,518 genes to orthogroups. Zea mays was excluded from this orthogroup run because a source-matched Zm00001d.1 protein FASTA was unavailable.

\suppmethodsection{Input inventory and integrity checks}

The governed manifest contained 121 audited standardized H5AD resources. After eligibility filtering, 44 unique H5AD objects contributed to discovery; their filenames, byte sizes and SHA-256 checksums are reported in Table S1b. PRJCA014789 was audited with the collection but excluded from discovery because it provided only one eligible vascular target label, and was reserved before scoring as the Populus phloem test. For Table S2, each readable object was opened with anndata in read-only backed mode and only obs metadata were accessed. The standardized source-label column cell\_type\_original was preferred, with celltype\_after permitted only as a recorded fallback. The scan produced 924 unique dataset--source-label pairs representing 212 distinct source-label strings and 3120007 cells. The corresponding manifest total was 3120007 cells.

Dataset identity, accession, species, tissue, matrix storage, and source-annotation checksums were taken from dataset\_manifest.tsv and dataset\_source\_metadata.tsv. Source-label counts were calculated directly from the standardized annotation field. Expression values were not consulted. A dataset lacking a recognized source-label field would have been retained as a source\_unavailable row rather than omitted; the inventory records the number of such datasets.

\suppmethodsection{Label governance and hierarchical mapping}

The global ontology candidate inventory contains 212 distinct source-label strings and explicitly retains Unknown. Table S2 maps observed dataset-label pairs only through frozen governance assets. The panel label freeze contains 29 panel--label decisions, of which 27 rows are approved for the pilot hierarchy and two source labels are excluded. The accepted hierarchy maps fine source labels to coarse root, leaf, or vascular families. Many-to-one mappings preserve the original source label, ontology identifier, governed label, parent panel, mapping rule, approval scope, and cell count.

A dataset-label pair receives mapped status only when the fixed dataset-panel label-count record and accepted hierarchy support that panel-specific mapping. Unknown labels remain unknown. Explicit exclusions remain excluded. Multiple or ambiguous ontology candidates remain ambiguous. All other observed labels remain unmapped, even when they have an exact lexical ontology candidate, because an ontology suggestion is not equivalent to a governed analysis label. This rule prevents expression-based relabelling and preserves the complete unresolved vocabulary.

\suppmethodsection{Reference mapping, orthogroup construction, and bridges}

Source H5AD \texttt{var\_names} were mapped to species reference-gene candidates with the fixed \texttt{gene\_\allowbreak mapping\_\allowbreak candidates\_v1} bridge. The mapping assessment covered 81 datasets from eight species, contained 1,611,655 feature rows, mapped 1,587,599 candidate rows, and reported no ambiguous candidate rows, target collisions, or gene rows with multiple representative proteins. These mappings remained candidates pending acceptance review and were not interpreted as experimental orthology validation.

Representative proteins were grouped by the seven-species OrthoFinder run. The resulting \texttt{gene\_\allowbreak orthogroup\_\allowbreak candidates\_v1} bridge linked mapped reference genes to 18,383 candidate orthogroups. Dataset-level orthogroup coverage was calculated as orthogroup-assigned source features divided by total source features. Seventy-one of 81 evaluated datasets passed the provisional 60\% gate. Table S3 expands these dataset-level values across the three frozen panels, retaining blank mapping fields for manifest datasets without an evaluated orthogroup bridge.

\suppmethodsection{Rank encoding and feature construction}

The executed rank representation used nonnegative source expression values from the frozen input layer. Within each cell, nonzero features were ordered by decreasing value, with exact ties resolved by immutable source feature order. The top 128 features were assigned weights 1/log2(rank+1), multiple source features mapping to the same orthogroup were collapsed by their maximum weight, and the sparse orthogroup vector was L2-normalized per cell. This design reduced dependence on incomparable expression magnitudes while retaining within-cell feature ordering.

Marker-rank feature v2 was created for the approved leaf, root, and vascular panels in read-only backed mode. It contained 40,417 leaf cells, 74,821 root cells, and 29,851 vascular cells over the 18,383-orthogroup vocabulary. The 44 discovery objects contributed 65 object--panel combinations (23 root, 19 leaf and 23 vascular), all with nonzero mapped features. Rank features were intended for cross-dataset marker and programme analysis, not differential-expression inference.

\suppmethodsection{Programme discovery}

Within each species and panel, each accepted coarse class was contrasted against the other accepted coarse classes. Orthogroups with positive mean normalized-rank-score differences were sorted in descending order. The frozen discovery rule retained the top 100 orthogroups per eligible species--class comparison, required at least 100 class cells and 100 other cells, and nominated candidates recurring in at least three species.

The v2 discovery generated 4,200 species-level marker rows from 42 eligible comparisons and 102 recurrent candidates: 39 leaf, 28 root and 35 vascular. These candidates came from the observational discovery collection and were not independent validation.

The GSE268881 xylem programme was selected separately as the first 18 positive orthogroups in the Arabidopsis vascular-xylem ranking in \texttt{coarse\_marker\_orthogroups\_v2.tsv}. It was not the 18-member recurrent xylem set produced by the top-100/minimum-three-species rule. This top-18 cut-off was a fixed heuristic chosen before external scoring, not an outcome-optimized threshold. The programme source, Arabidopsis rank and score difference are recorded for every member in Table S4; six of 18 members were later rediscovered by the recurrence rule (Table S7c).

\suppmethodsection{Freeze protocol and leakage controls}

Labels, dataset panels, the hierarchy, programme membership, contrasts, scoring parameters, random seeds and decision rules were fixed before each corresponding validation run. Source objects were opened read-only, and versioned records retained input identities, protocol changes and superseded runs.

Internal leave-one-study-out analyses selected the dominant gene using all but one dataset within an orthogroup and species, then evaluated the selected gene in the omitted dataset. This reduced same-study selection bias but remained internal evidence. External validation programmes were pre-discovered, and target contrasts were frozen before scoring. Missing mappings and excluded labels were not replaced with related genes or labels.

\suppmethodsection{Validation rounds}

Study-level robustness was assessed for 19 priority orthogroups by leaving out each available study within a species. The resulting family statuses summarize direction and dominant-gene stability across held-out studies. Clean supported trees were then used for descriptive dominant-gene compactness analyses. Neither different dominant genes nor dispersed tree positions were interpreted as paralog substitution, and the final analysis found no family satisfying the strict joint expression-plus-supported-subclade criterion.

Independent rice validation v2.1 used the RNA counts layer from GSE232863, encoded 49,273 selected cells, mapped 14,136 orthogroups, used top $K=128$, and compared each fixed programme with 500 random programmes under seed 20260914. Version 2.1 corrected only the per-contrast sample-direction filter; programmes, contrasts, thresholds and seed were unchanged. Three of five primary contrasts passed. Replicate-level raw rows were unavailable and were not reconstructed.

External validation v3 fixed the rank encoding, mapped-member programme score, single-gene and expanded-family subsets, off-target controls, 500 random programmes and replicate-direction rule before expression scoring. The tables retain the negative Populus phloem result and distinguish unavailable values from zero. GSE268881 Esa and Sir provide held-out agreement with label-transfer-derived source assignments and within-study repeatability, not independent biological validation or replication across studies; Spa remained ineligible under the metadata gate.

\suppmethodsection{Nested labels, marker-overlap exclusion, and parameter robustness}

The nested-label analysis used the GSE268881 RNA counts layer and three fixed endpoints: xylem versus all other root labels; xylem versus pericycle, phloem and procambium; and stele/vascular labels versus non-stele root labels. The same 18-orthogroup score was evaluated within the same cohorts, species and biological samples, but each endpoint necessarily contained a different set of positive and negative cells. Results were calculated for pooled cells and separately for R1 and R2 in Esa and Sir. The $K=128$ source-order condition reproduced the prior pooled and replicate AUROCs to $10^{-12}$.

Parameter robustness covered K values 64, 128, and 256 and the immutable source order plus ten deterministic global feature-order permutations. Across the mapped-member programme, all 18 species--endpoint--scope groups retained AUROC above 0.5 and positive standardized mean difference in all 33 conditions per group. These direction checks are descriptive robustness results rather than new confirmatory endpoints.

Marker-overlap sensitivity removed five of the 18 fixed orthogroups documented as overlapping published marker evidence. The complete label-assignment or label-transfer feature set was unavailable. Therefore the analysis supports robustness to removal of documented overlap only and does not completely exclude label circularity. Integrated scale.data was not used as a substitute marker set.

\suppmethodsection{Discovery-threshold and clade-balance sensitivity}

The discovery sensitivity analysis reconstructed the top-100/minimum-three-species candidate set and then evaluated top-$N$ values of 25, 50 and 100 and recurrence thresholds of two to five species. At a minimum of three species, 16, 35 and 102 candidates were retained for top 25, 50 and 100, respectively.

Cross-clade presence required at least one eudicot and one monocot, whereas the strict balance descriptor required at least two species from each clade. Among the frozen 102 candidates, 62 had cross-clade presence and seven met the strict two-plus-two descriptor. Leaf collections were eudicot-only and thus could not pass a cross-clade rule; this was recorded as a design limitation rather than negative evidence. Six of the submitted 18 xylem-programme orthogroups were independently rediscovered by the frozen vascular-xylem top-100/minimum-three rule.

\suppmethodsection{Statistical summaries}

Cell-level discrimination was summarized by AUROC, average precision, and standardized mean difference. Average precision was interpreted relative to class prevalence. Replicate scopes were reported separately from pooled cells. Random-programme comparisons used the frozen empirical probability (1 + number of random AUROCs at least as large as the target AUROC)/(number of random programmes + 1). Direction rules used AUROC greater than 0.5 and positive standardized mean difference; these rules were frozen and were not optimized after outcome inspection.

No missing metric was converted to zero. Counts and mapping fractions in Table S3 are numeric; each reported fraction was recomputed as mapped\_features/total\_features, checked to lie within [0,1] and compared with the recorded value to $10^{-12}$. Descriptive sensitivity analyses did not create new hypothesis-testing thresholds.

\suppmethodsection{Reproducibility and data integrity}

Input and output identities were recorded with SHA-256 hashes in the reproducibility repository. Supplementary Table S1 states the scope of every checksum: direct source-object hashes are reported only when available, while metadata or derived-record hashes are labelled as such. Table S1b supplies filenames, byte sizes and direct SHA-256 checksums for all 44 standardized H5AD objects used in discovery.

Tables were checked for unique record keys, nonnegative cell and feature counts, valid fraction bounds, numerator/denominator agreement and preservation of missing values. Scanned cell counts were compared with the manifest, and all observed label strings were retained.

\suppmethodsection{Known gaps}

Dataset-specific source-label field names before standardization are unavailable for many annotations; Table S2 therefore reports the standardized H5AD field and available source-annotation hash. Only 29 panel--label decisions and 27 accepted hierarchy rows were available, so other observed labels remain unknown, excluded, ambiguous or unmapped. Orthogroup coverage is available for 81 of 121 discovery datasets; the remaining 40 retain blank numerators and fractions. Gene and orthogroup mappings are candidate bridges rather than experimentally validated one-to-one orthologues. Functional conservation, causality, complete absence of label circularity and paralogue substitution remain unestablished.

